% =============================================================================
% Fragmento para integracao ao relatorio (projeto Riachuelo, PRO3474)
% -----------------------------------------------------------------------------
% Ponto de integracao sugerido: dentro da secao "Refinar o ciclo de vida do
% produto e clientes do projeto" (\label{sec:refinar_ciclo_de_vida}), como
% subsecao posicionada apos "Definicao dos clientes do projeto"
% (\label{subsec:definicao_dos_clientes}).
%
% Versao revisada (stop-slop): corta contrastes "nao X, e Y" repetidos, troca
% vozes passivas por sujeito ativo, remove adverbios e enfase redundante.
%
% Dependencia bibliografica: \cite{rozenfeld2006gestao} ja e usada no
% documento principal.
% =============================================================================

\subsection{Segmentação de clientes e usuários prioritários}
\label{subsec:segmentacao_clientes_usuarios}

\subsubsection{Estrutura de stakeholders do projeto}
\label{subsubsec:segmentacao_estrutura}

A equipe define quatro clusters de stakeholders: Cliente Premium, Cliente Popular, Operação da Riachuelo e Gestão da Riachuelo. Cliente Premium e Cliente Popular são consumidores finais das lojas: usam o produto na ponta da jornada de compra e correspondem, na classificação de Rozenfeld et al. \cite{rozenfeld2006gestao}, a clientes externos. Operação da Riachuelo e Gestão da Riachuelo são clientes internos: executam ou decidem sobre o sistema que viabiliza a solução, sem comprar roupa na loja.

\textbf{Cliente Premium.} Consumidor de lojas ou formatos voltados a público de maior poder aquisitivo. Interage com a jornada completa em estudo: escolhe o produto, paga e sairia da loja sem passar pelo caixa, caso a solução avance. A observação de campo aponta expectativas diferentes de autonomia, atendimento e tecnologia em relação ao público popular. Essa diferença justifica o cluster, e a pesquisa desta etapa vai testá-la em campo.

\textbf{Cliente Popular.} Consumidor do varejo de massa. Também é usuário final da jornada, mas parte de um padrão de compra observado como mais assistido pelo vendedor, inclusive na oferta de crédito. A hipótese de comportamento distinto do Cliente Premium exige um contraponto real para ser testada.

\textbf{Operação da Riachuelo.} Reúne vendas, estoque, reposição, caixa, segurança e demais funções que tocam a loja no dia a dia. Manuseia, monitora ou sofre o impacto direto de qualquer mudança de processo físico na loja, efeito que não aparece na voz do consumidor. Esse efeito justifica tratar esse público como cluster próprio.

\textbf{Gestão da Riachuelo.} Reúne middle management e alta gestão. Avalia a solução por indicadores, resultado operacional, retorno econômico, risco e impacto estratégico. Em um modelo B2B, este grupo aprova a compra da solução; quem usa o produto na loja não participa dessa decisão.

\subsubsection{Escopo desta etapa de pesquisa}
\label{subsubsec:segmentacao_escopo}

Os quatro clusters formam a estrutura global de stakeholders do projeto. Esta etapa de pesquisa investiga apenas Cliente Premium e Cliente Popular, porque a equipe ainda não tem acesso estruturado à operação interna da Riachuelo para conduzir entrevistas ou surveys consistentes com Operação e Gestão. Este documento evita listar hipóteses de necessidades para esses dois clusters: inventar necessidades sem base em campo comprometeria a etapa seguinte de investigação. Operação e Gestão permanecem registrados como stakeholders relevantes; a pesquisa com eles ocorre em etapa posterior, quando esse acesso existir.

\subsubsection{Diferenças relevantes entre Cliente Premium e Cliente Popular}
\label{subsubsec:segmentacao_diferencas}

A Tabela~\ref{tab:segmentacao_grupos} resume o que já foi observado em campo e o que ainda é hipótese a testar nesta etapa.

\begin{table}[h!]
    \centering
    \caption{Cliente Premium e Cliente Popular: observações e hipóteses}
    \begin{tabular}{p{2.6cm}p{4.6cm}p{4.6cm}}
        \toprule
        \textbf{Dimensão} & \textbf{Cliente Premium} & \textbf{Cliente Popular}\\
        \midrule
        Papel do vendedor (observação de campo) & Menor influência sobre a decisão de cartão de crédito & Maior influência sobre a decisão de cartão de crédito\\
        Autoatendimento (observação de campo) & Maior preferência declarada & Menor preferência declarada\\
        Dependência do vendedor (hipótese a testar) & Pode estar diminuindo, com valorização crescente de autonomia & Padrão de dependência ainda não medido, pode seguir tendência semelhante em outro ritmo\\
        Lógica de exposição de produto (observação de campo) & Não caracterizada nesta etapa & Prefere ver muitas peças expostas (lógica de abundância)\\
        \bottomrule
    \end{tabular}
    \label{tab:segmentacao_grupos}
\end{table}

A preferência entre autonomia e atendimento assistido varia por momento da jornada e por cluster. A pesquisa precisa localizar em que momentos o consumidor valoriza o atendimento humano, em que momentos esse atendimento vira necessidade e em que momentos o consumidor prefere resolver sozinho. Tratar autonomia como vantagem automática arrisca construir uma solução que remove suporte no momento em que o cliente ainda precisa dele.

\subsubsection{Hipóteses a validar nesta etapa}
\label{subsubsec:segmentacao_hipoteses}

\begin{itemize}
    \item \textbf{H1}: no Cliente Premium, a dependência do vendedor para a decisão de compra está diminuindo, com valorização crescente de autonomia durante a jornada.
    \item \textbf{H2}: o mesmo movimento pode ocorrer no Cliente Popular, em velocidade ou intensidade diferente.
    \item \textbf{H3}: a fricção mais relevante da jornada de compra pode incluir disponibilidade de produto, dificuldade de encontrar itens na loja ou integração entre canal físico e digital, além de fila e pagamento.
    \item \textbf{H4}: os dois clusters valorizam atendimento humano em momentos distintos da jornada, como escolha e prova de roupa em contraste com pagamento; essa valorização muda ao longo do percurso de compra.
\end{itemize}

O survey e o focus group descritos nas seções seguintes deste relatório decidem se cada hipótese se confirma.

\subsubsection{Implicações para a pesquisa}
\label{subsubsec:segmentacao_implicacoes}

A pesquisa busca se aprofundar do comportamento real do consumidor. O survey deverá medir a relevância e a frequência das fricções e comparar Cliente Premium e Cliente Popular por meio de afirmações avaliáveis em escala, sem sugerir a ideia final. O focus group aprofundará por que essas fricções ocorrem, em que contexto e o que o consumidor faz hoje para contorná-las. O focus group apresentará a trava antifurto conectada ao aplicativo apenas na etapa final, depois de mapear a experiência atual.
